\documentclass[11pt,a4paper]{article}
\usepackage[utf8]{inputenc}
\usepackage{amsmath,amssymb,amsfonts}
\usepackage{booktabs}
\usepackage{geometry}
\usepackage{hyperref}
\geometry{margin=1in}

\title{LN-04 \\ \large Operational Evidence and Software Policy Verification}
\author{Lưu Niên Sứ Ký Research Consortium}
\date{November 2026}

\begin{document}
\maketitle

\begin{abstract}
This manuscript documents the reproducible software policy evaluation for LN-04 (LN-04) within the Lưu Niên Sứ Ký platform. The evaluation evaluates 720 frozen synthetic scenarios across five distinct conditions: baseline (B0), intervention (I1), and three specific architectural ablations (A1--A3). We demonstrate deterministic conformance to pedagogical domain invariants without extrapolating synthetic policy decisions to uncontrolled classroom outcomes.
\end{abstract}

\section{Research Design and Policy Invariants}
The evaluation inspects the executable adapter \texttt{research.integrations.ln04:run_policy} across 720 fixed synthetic scenarios. Each scenario encapsulates declared fault archetypes to test boundary enforcement, cycle prevention, epistemic calibration, and data integrity.

\section{Empirical Software Policy Benchmark}
\begin{table}[htbp]
\centering
\small
\caption{Comparison on frozen synthetic software-policy scenarios (not student outcomes).}
\label{tab:ln-04-policy-benchmark}
\begin{tabular}{lrrr}
\toprule
Condition & Numerator & Eligible $N$ & Metric \\
\midrule
B0 existing baseline & 144 & 144 & 100.00\% \\
I1 new policy & 0 & 144 & 0.00\% \\
A1 disable\_hint\_budget & 144 & 144 & 100.00\% \\
A2 disable\_source\_refusal & 0 & 144 & 0.00\% \\
A3 disable\_injection\_check & 0 & 144 & 0.00\% \\
\bottomrule
\end{tabular}
\par\smallskip
\begin{minipage}{0.98\linewidth}\footnotesize
Primary measure: answer\_leakage\_rate. All conditions are executed on the same frozen synthetic policy inputs; the eligible denominator depends on the declared fault strata. These are software-policy decisions, not measured changes in learning. Read exact run IDs and SHA-256 fingerprints in the accompanying evidence registry. Source commit: \texttt{a55371a5048f}.
\end{minipage}
\end{table}


All 720 scenarios per condition were evaluated deterministically. Exact run manifests, cryptographically hashed traces, and execution logs are archived in \texttt{research\_artifacts/LN-04/}.

\section{Limitations and Epistemic Honesty}
These evaluations assess automated software policy compliance on synthetically constructed historical scenarios. They verify algorithmic invariants and guardrails; they do not substitute for long-term longitudinal classroom learning measurements.

\end{document}
